\cotefile{75}
\scanpage{26}
c) Orientation d'un ensemble fini $I$\inserted{(cardinal $n$).} La donnée de $I$ équivaut à celle d'un torseur
\[
 P=\operatorname{Isom}(\deleted{[1,n]_{\unsure{\mathbb Z}}}I_n,I)=\operatorname{numérot}(I)
\]
$(I_n\mathrel{\unsure{:|=}}\{m\in\mathbb N^{*}\mid m\leq n\}=\{\deleted{0}1,2,\ldots,n\})$ sous le groupe
\[
 \deleted{\unsure{\mathfrak S|S}}\mathfrak S_n=\mathfrak S_{I_n}=\operatorname{Aut}(I_n).
\]

On a un hom canonique ``signature''
\diagram{\begin{tikzcd}[ampersand replacement=\&]
 \mathfrak S_n \arrow[r,"s=\mathrm{sgn}"] \& \{\pm1\} \& \text{noyau }\mathfrak A_n
\end{tikzcd}}
et par définition l'ens des 2 orientations de $I$ est
\[
 \unsure{\omega|w}_I=\deleted{\operatorname{Isom}(}P\times_{\mathfrak S_n}\{\pm1\}\simeq P/\mathfrak A_n.
\]

\uline{Compatibilités avec} a) \uline{et} b)

\uline{Avec} b) On a sur $E_{(I,k)}=k^{I}$ un produit scalaire canonique
\[
 (x,y)=\sum_{i\in I}x_i y_i
\]
pour lequel $(e_i)$ est une base orthonormale. Ainsi $\deleted{E}I$ \deleted{\unsure{donne}}\inserted{\unsure{donne}} naissance à un module quadratique $E(I,\unsure{x|k|X})$. Ceci posé, on \inserted{définit} une bijection canonique
\[
 \omega_{E(I,x)}\simeq\omega_I
\]
\margin{Si $i_1,\,\cdot\cdot\,i_n$ est une numérotation de $I$, $e_{i_1}\,\cdot\cdot\,e_{i_{\unsure{n}}}$ est une base \inserted{\unsure{ord.}} orthon. de $E$, \ldots}
par transitivité dans
\diagram{\begin{tikzcd}[ampersand replacement=\&,column sep=large]
 \mathfrak S_n \arrow[r,hook,"\mathrm{incl.}"] \arrow[rr,bend right=24,"\mathrm{sgn}"] \& \mathrm{O}(\unsure{N|n},k) \arrow[r,"\det"] \& \{\pm1\}
\end{tikzcd}}

En particulier, on trouve (comp. avec a).),
% The bottom annotation is written on two lines, with its second line at the cut edge.
% The arrow leaves the closing parenthesis of the low E(I,R), at the southeast of the first cell.
\diagram{\begin{tikzcd}[ampersand replacement=\&,column sep=large,row sep=small]
 \omega_{E(I,\mathbb R)} \arrow[dr,start anchor=south east,end anchor=west] \& \simeq\omega_I \\
 \& \begin{gathered}\text{considéré}\\[-2pt]\text{comme \unsure{variété} \unsure{via|vie} \unsure{se pense|ses pts|ses points}}\end{gathered}
\end{tikzcd}}
% archive sheet number: 25
